\section{Introdução}
\label{cap:introducao}

Este relatório apresenta os resultados do primeiro Trabalho Computacional da disciplina de Reconhecimento de Padrões. O objetivo principal é implementar e avaliar o desempenho de diferentes algoritmos de classificação supervisionada: K-Vizinhos Mais Próximos (\textit{K-Nearest Neighbors} - KNN), Classificador de Distância Mínima ao Centróide (\textit{Minimum Distance to Centroid} - MDC) e o Classificador de Máxima Correlação de Cosseno (\textit{Maximum Cosine Correlation} - MAXCO).

A análise é conduzida sobre o conjunto de dados \textit{Vertebral Column}, que contém medições biomecânicas da coluna vertebral de pacientes, classificados em três categorias: Hérnia de Disco, Normal e Espondilolistese. O trabalho explora diferentes métricas de distância para o KNN (através da distância de Minkowski com variados ordens), versões robustas do MDC utilizando a mediana, e a influência da proporção de dados de treinamento no desempenho final dos modelos.

A estrutura deste relatório compreende a descrição da metodologia empregada, incluindo o pré-processamento e o protocolo experimental, seguida da apresentação e discussão dos resultados obtidos através de múltiplas rodadas de treinamento e teste.
